% Section 6: conclusions. Numbers: generated/numbers.tex.
\section{Conclusions}\label{sec:conclusions}
The vertical support stiffness of a wheelset bearing can be identified from in-service axle-box
acceleration without measuring the rail roughness, by pooling passes at different speeds and payloads
in a debiased Whittle likelihood in which the roughness levels are free nuisance parameters. For a
representative freight wagon on ballasted track, simulations show that:
\begin{itemize}
\item with 50 passes at 60--110~km/h the support stiffness is estimated with a bias of
  \HLTwentyBias--\HLFortyBias\% and a root-mean-square error of \HLTwentyRMSE--\HLFortyRMSE\%, close to
  the Cram\'er--Rao bound at 20~dB and \HLFortyRatioCRB{} times it at 40~dB, where intervals are
  calibrated by simulation;
\item the speed spread is the design variable that makes the support stiffness identifiable, and the
  wheel-to-wheel coupling through the rail must be modelled, since ignoring it biases the estimate by
  \IgnoreCouplingForty\%;
\item the absolute level carries a model-form uncertainty of about \BudgetCoest\%, so the method is a
  trend monitor: re-estimating the foundation and contact stiffness in each period, a decrease of
  \MDCCalFiftyTwenty--\MDCCalFiftyForty\% is detectable with 50 passes per period, and a 20\% change of
  the foundation raises \DriftMCCoestFA{} false alarm in \DriftMCCoestN{} trials;
\item a loss of housing, adapter or pad stiffness reaches the identified stiffness directly, while
  clearance growth is diluted by the stiff roller contacts, so the method reports support-stiffness
  loss without attributing it to a mechanism.
\end{itemize}
The next steps are a field test with a controlled change of support stiffness; application to
freight wagons of Indian Railways, once vehicle-specific values are available from RDSO; an extension of
the model to the pinned--pinned frequency band; and calibrated intervals for the track-side use of the
same likelihood.
